% SECCIÓN 6: ANÁLISIS DE LA JURISPRUDENCIA PERUANA — PATRONES Y CRITERIOS DOMINANTES
%
% Objetivos de esta sección:
%   1. Presentar el corpus de sentencias analizadas (Corte Suprema + Cortes Superiores,
%      período 2010–2024) y los criterios de selección.
%   2. Identificar los patrones jurisprudenciales en relación con:
%      a) Lex artis: ¿cómo la determinan los jueces? ¿recurren a peritos? ¿a protocolos?
%      b) Consentimiento informado: ¿cuándo se considera insuficiente? ¿qué consecuencias?
%      c) Daño médico: ¿qué tipos de daño se reconocen? ¿cómo se cuantifican?
%   3. Analizar la coherencia/incoherencia entre los criterios de la Sala Civil
%      Permanente y la Sala Civil Transitoria de la Corte Suprema.
%   4. Identificar los criterios determinantes del resultado (fundada/infundada):
%      ¿qué elemento es más frecuentemente deficiente — lex artis, nexo causal o daño?
%   5. Señalar las lagunas jurisprudenciales y propuestas de lege ferenda.
%
% ESTRUCTURA RECOMENDADA:
%   \subsection{Corpus y metodología}
%   \subsection{Criterios sobre la lex artis}
%   \subsection{Criterios sobre el consentimiento informado}
%   \subsection{Criterios sobre el daño médico y su valoración}
%   \subsection{Patrones sistémicos y vacíos jurisprudenciales}
%
% Extensión estimada: 2.500–3.500 palabras (sección central del artículo)

\section{Análisis de la jurisprudencia peruana: patrones y criterios dominantes}
\label{sec:jurisprudencia}

